\setchoices{label=(\Alph*)}

\begin{question}[index=1]
当 $x\to0^+$ 时，下列式子正确的个数为
\par\circlednumber{1} $\dfrac{1}{\ln x}\asympteq\dfrac{1}{\ln\sin x}$；
\par\circlednumber{2} $\displaystyle\int_0^x t^2\dd t\asympteq\int_0^x\sin t^2\dd t$；
\par\circlednumber{3} $(1+\sin x)^x-1\asympteq(1+x)^{\sin x}-1$。
\begin{choices}[columns=4]
\item $0$
\item $1$
\item $2$
\item $3$
\end{choices}
\end{question}
\QuestionBreak

\begin{question}[index=2]
设曲面 $\Sigma:z=x^2+y^2$（$0\leq z\leq1$），取下侧，则
\[
\iint_{\Sigma}(x+y)\dd y\dd z+(z+2xy+2x^2)\dd x\dd y=
\]
\begin{choices}[columns=4]
\item $-\dfrac{\pi}{2}$
\item $-\dfrac{\pi}{4}$
\item $\dfrac{\pi}{4}$
\item $\dfrac{\pi}{2}$
\end{choices}
\end{question}
\QuestionBreak

\begin{question}[index=3]
过点 $(1,2,0)$ 的直线 $L$ 与直线 $L_1:\dfrac{x-2}{1}=\dfrac{y}{2}=\dfrac{z-6}{-2}$ 相交，且平行于曲面 $\Sigma:z=x^2+2y^2$ 在点 $(1,1,3)$ 处的切平面，则直线 $L$ 的方程为
\begin{choices}[columns=2]
\item $\dfrac{x-1}{2}=\dfrac{y-2}{0}=\dfrac z4$
\item $\dfrac{x-1}{2}=\dfrac{y-2}{-1}=\dfrac z0$
\item $\dfrac{x-1}{0}=\dfrac{y-2}{1}=\dfrac z4$
\item $\dfrac{x-1}{1}=\dfrac{y-2}{0}=\dfrac z1$
\end{choices}
\end{question}
\QuestionBreak

\begin{question}[index=4]
设 $f(x)$ 在 $(0,+\infty)$ 内二阶可导，则下列说法正确的是
\begin{choices}[columns=1]
\item 若 $f(x)$ 在 $(0,+\infty)$ 内有界，则 $\displaystyle\lim_{x\to+\infty}f'(x)=0$
\item 若 $f(x)$ 和 $f'(x)$ 在 $(0,+\infty)$ 内均有界，则 $\displaystyle\lim_{x\to+\infty}f''(x)=0$
\item 若 $f'(x)$ 在 $(0,+\infty)$ 内有界，则 $f(x)$ 在 $(0,+\infty)$ 内一定有界
\item 若 $f(x),f''(x)$ 在 $(0,+\infty)$ 内有界，则 $f'(x)$ 在 $(0,+\infty)$ 内一定有界
\end{choices}
\end{question}
\QuestionBreak

\begin{question}[index=5]
设 $\vect A=\mat{\dfrac32&-\dfrac12\\[3pt]\dfrac12&\dfrac12}$，向量 $\vect\alpha$ 满足 $\vect A\vect\alpha=\mat{1\\1}+\vect\alpha$，则 $\vect A^{2025}\vect\alpha$ 不可能为
\begin{choices}[columns=4]
\item $\mat{2025\\2023}$
\item $\mat{2026\\2024}$
\item $\mat{2027\\2025}$
\item $\mat{2027\\2026}$
\end{choices}
\end{question}
\QuestionBreak

\begin{question}[index=6]
设 $3$ 阶实对称矩阵 $\vect A=(\vect\alpha_1,\vect\alpha_2,\vect\alpha_3)$，且 $\vect\alpha_1+\vect\alpha_2+\vect\alpha_3=(1,1,1)^{\mathrm T}$，又 $a\in\mathbb R$ 使得
\[
\vect A^3+(a^5-1)\vect A^2+2a^3\vect A+a\vect E=\vect O,\quad \operatorname{tr}(\vect A)=1,
\]
则下列结论中
\par\circlednumber{1} $r(\vect A)=1$；
\par\circlednumber{2} $\vect A$ 的特征值为 $0,0,1$；
\par\circlednumber{3} $\vect A$ 与矩阵 $\mat{1&0&0\\0&2&0\\0&0&0}$ 合同；
\par\circlednumber{4} $\mat{1\\1\\-2}$ 为 $\vect A\vect x=\vect0$ 的解。
\par 正确的个数为
\begin{choices}[columns=4]
\item $1$
\item $2$
\item $3$
\item $4$
\end{choices}
\end{question}
\QuestionBreak

\begin{question}[index=7]
矩阵 $\vect P$ 满足 $\vect P\vect P^{\mathrm T}=\vect E$，下列四个式子中必成立有几个？
\par\circlednumber{1} $\vect P$ 的行向量组为标准正交向量组；
\par\circlednumber{2} $\vect P$ 的列向量组为标准正交向量组；
\par\circlednumber{3} $\vect P$ 为正交矩阵；
\par\circlednumber{4} $\vect P$ 为可逆矩阵。
\begin{choices}[columns=4]
\item $1$
\item $2$
\item $3$
\item $4$
\end{choices}
\end{question}
\QuestionBreak

\begin{question}[index=8]
设随机变量 $X$ 与 $Y$ 独立同分布，且 $P\{X=k\}=\left(\dfrac12\right)^k,\ k=1,2,\dots$，$Z=\max\{X,Y\}$，则 $P\{Z=3\}=$
\begin{choices}[columns=4]
\item $\dfrac38$
\item $\dfrac7{32}$
\item $\dfrac{11}{64}$
\item $\dfrac{13}{64}$
\end{choices}
\end{question}
\QuestionBreak

\begin{question}[index=9]
设 $X_1,X_2,\dots,X_n,X_{n+1}$ 为来自总体 $X\dist N(\mu,\sigma^2)$ 的简单随机样本，$\displaystyle\overline X=\dfrac1n\sum_{i=1}^nX_i$，$T=a(X_{n+1}-\overline X)^2$，若使 $T\dist\chi^2(1)$，则 $a=$
\begin{choices}[columns=4]
\item $\dfrac{n}{(n+1)\sigma^2}$
\item $\dfrac{n\sigma^2}{n+1}$
\item $\dfrac{n}{(n+1)\sigma}$
\item $\dfrac{n+1}{n\sigma}$
\end{choices}
\end{question}
\QuestionBreak

\begin{question}[index=10]
设随机事件 $A$ 与 $B$ 互不相容，$B$ 与 $C$ 相互独立，且 $P(A)=P(B)=0.4,\ P(C)=0.5$，则 $P(C-A\mid A\cup BC)=$
\begin{choices}[columns=4]
\item $\dfrac15$
\item $\dfrac16$
\item $\dfrac13$
\item $\dfrac14$
\end{choices}
\end{question}
